\documentclass[11pt,a4paper]{article}

\usepackage[margin=1in]{geometry}
\usepackage{amsmath,amssymb,amsthm}
\usepackage{hyperref}
\usepackage{url}
\usepackage{booktabs}
\usepackage{longtable}
\usepackage{enumitem}
\usepackage{xcolor}

\hypersetup{colorlinks=true, linkcolor=blue!50!black, citecolor=blue!50!black,
            urlcolor=blue!50!black}
\setlength{\emergencystretch}{1.5em}

\newtheorem{definition}{Definition}
\newtheorem{conjecture}{Conjecture}
\newtheorem{lemma}{Lemma}
\newtheorem{proposition}{Proposition}
\newtheorem{observation}{Observation}
\newtheorem{remark}{Remark}
\newtheorem{question}{Question}

\newcommand{\Q}{\mathcal{Q}}
\newcommand{\Pool}{\mathcal{P}}
\newcommand{\R}{\mathcal{R}}
\newcommand{\Rflat}{\mathcal{R}_{\mathrm{flat}}}
\newcommand{\Radd}{\mathcal{R}_{+}}
\newcommand{\code}[1]{\texttt{#1}}

\title{Ladhe's Quad Conjecture:\\
Every Prime Quadruplet Is an Additive/Multiplicative/Exponential\\
Combination of Its Predecessors\thanks{Correspondence:
\texttt{spalgorithm@gmail.com}. Paper and dataset are archived on Zenodo,
\url{https://doi.org/10.5281/zenodo.22286847}, and released under
CC~BY~4.0; the code is released under the licence of the repository
\url{https://github.com/SPAlgorithm/LE}.}}

\author{Shubham Ladhe \and Saloni Ladhe \and Pankaj Ladhe}

\date{September 2026\\[2pt] {\small DOI: \href{https://doi.org/10.5281/zenodo.22286847}{10.5281/zenodo.22286847}}}

\begin{document}
\maketitle

\begin{abstract}
Prime quadruplets are the constellations $(p,\,p+2,\,p+6,\,p+8)$ of four
primes, the densest possible packing of four primes beyond $(5,7,11,13)$.
We order them into a chain $\Q_1=(11,13,17,19)$, $\Q_2=(101,103,107,109)$,
$\Q_3=(191,193,197,199)$, \dots, seeded by the \emph{royal quad}
$\Q_0=(2,3,5,7)$, and observe that every member of every quadruplet can be
written as the largest member of the previous quadruplet plus a sum of
distinct members of earlier quadruplets, for example
$13001 = 9439 + 3461 + 101$ and
$35545421 = 35524459 + 18049 + 2089 + 821 + 3$.
The precise form of the observation is that the gap $D$ between a
quadruplet member and its base is a sum of at most $8$ distinct earlier
members, with exactly ten exceptional gaps, all at most $382$, that need a
small expression in the royal members $2,3,5,7$; allowing one such
expression in every gap brings the bound down to six terms with no
exception at all. We call this \emph{Ladhe's Quad Conjecture} and verify it
for all $28{,}387$ prime quadruplets below $10^9$, i.e.\ for $113{,}544$
quadruplet members, with an independent checker. Two further
representations of each gap, one that prefers multiplication and one that
prefers powers of chain members, are also computed and released. We also
prove, from a two-line lemma and computed base cases, that the chain is an
additive basis for all integers: every integer up to the next quadruplet
(indeed up to a bound that grows like the sum of all earlier members) is a
sum of distinct chain members plus one expression in $2,3,5,7$ with
$+,-,\times$; this is unconditional below $10^9$ and rests beyond that
only on a mild growth condition. We explain carefully what the conjecture
does and does not say about locating the next quadruplet, give a heuristic
for why such representations are plentiful, and state open questions. The utility that builds the chain and
the full dataset are available at
\url{https://github.com/SPAlgorithm/LE}.
\end{abstract}

\section{Introduction}

A \emph{prime quadruplet} is a quadruple of primes of the form
\[
(p,\; p+2,\; p+6,\; p+8),
\]
the tightest admissible pattern of four primes. Apart from $(5,7,11,13)$,
every prime quadruplet has $p\equiv 11 \pmod{30}$, so its members end in the
digits $1,3,7,9$ in that order. The sequence of first members
$11, 101, 191, 821, 1481, 1871, 2081, \dots$ is OEIS A007530
\cite{OEIS_A007530}. Whether there are infinitely many prime quadruplets is
open; it is the case $k=4$ of the Hardy--Littlewood prime $k$-tuple
conjecture \cite{HardyLittlewood1923}, and the Bateman--Horn heuristic
\cite{BatemanHorn1962} predicts that the number of quadruplets up to $x$ is
asymptotically $C_4\, x/(\log x)^4$ with $C_4 \approx 4.15$. The strongest
unconditional results in this direction are the bounded-gap theorems of
Zhang \cite{Zhang2014} and Maynard \cite{Maynard2015}, which give
infinitely many pairs (and $m$-tuples) of primes within a bounded distance
but not the specific pattern $(0,2,6,8)$. Large quadruplets and related
clusters have been searched for computationally for decades
\cite{Forbes1999}.

This note records an additive structure that runs through the whole
sequence of prime quadruplets, at least as far as we have looked. Order the
quadruplets by size and call the $k$-th one $\Q_k$; put the four primes
$2,3,5,7$ in front as $\Q_0$, the \emph{royal quad}. Then every member of
$\Q_k$ is obtained from the largest member of $\Q_{k-1}$ by adding a few
distinct members of earlier quadruplets. Section~\ref{sec:conj} states this
precisely in three forms of increasing strength, Section~\ref{sec:cover}
proves a companion statement, that with subtraction allowed among the
royal members every integer below the next quadruplet is such a sum,
Section~\ref{sec:examples} walks through examples that can be checked by
hand, Section~\ref{sec:evidence} reports the computation up to $10^9$, and
Section~\ref{sec:next} discusses, with some care, what the statement means
for finding the next quadruplet.

The observation grew out of the authors' earlier work on additive
decompositions of primes, \emph{Ladhe's Conjecture} \cite{LadheConj}, which
concerns single primes written as sums of three primes and underlies a
hash-based signature scheme \cite{LadheScheme}. The present note is about
prime \emph{quadruplets} and is self-contained.

\section{Definitions and rules}
\label{sec:defs}

\begin{definition}[Quads]
$\Q_0=(2,3,5,7)$ is the \emph{royal quad}. For $k\ge 1$, $\Q_k$ is the
$k$-th prime quadruplet $(p,p+2,p+6,p+8)$ with $p\ge 11$, in increasing
order of $p$. We write its members as
$q_k^{(1)}=p,\; q_k^{(3)}=p+2,\; q_k^{(7)}=p+6,\; q_k^{(9)}=p+8$, the
superscript being the last decimal digit. The quadruplet $(5,7,11,13)$ is
not part of the chain, because $5$ and $7$ already belong to $\Q_0$, and
$\Q_0$ itself is not of the form $(p,p+2,p+6,p+8)$.
\end{definition}

\begin{definition}[Base, gap, pool]
For $k\ge 2$ the \emph{base} of $\Q_k$ is $m_k=q_{k-1}^{(9)}$, the largest
member of the previous quadruplet. For a member $t$ of $\Q_k$ the
\emph{gap} is $D=t-m_k>0$. The \emph{pool} available to $\Q_k$ is
$\Pool_k=\Q_1\cup\dots\cup\Q_{k-1}$, the set of all $4(k-1)$ primes of the
earlier quadruplets.
\end{definition}

\begin{definition}[Royal expressions]
A \emph{royal expression} is an arithmetic expression built from the
members $2,3,5,7$ of $\Q_0$, each used at most once, with the operations
$+$ and $\times$ only. Let $\R$ be the set of its values; $\R$ has $68$
elements between $2$ and $210$. Two subsets matter below: $\Radd$, the
$11$ values reachable by addition alone,
$\{2,3,5,7,8,9,10,12,14,15,17\}$, and $\Rflat$, the $38$ values with a
\emph{flat} expression, one in which no product contains a sum, such as
$(5\times 7)+(2\times 3)=41$ but not $5\times(7+(2\times3))=65$.
\end{definition}

\begin{definition}[Additive representation]
An \emph{additive representation} of a member $t$ of $\Q_k$ ($k\ge2$) is an
identity
\[
t \;=\; m_k \;+\; s_1+s_2+\dots+s_j \;+\; r,
\qquad s_1>s_2>\dots>s_j,\;\; s_i\in\Pool_k,\;\; r\in\R\cup\{0\},
\]
with all $s_i$ distinct and none equal to $m_k$. Its \emph{length} is
$j$ if $r=0$ and $j+1$ otherwise. Several $s_i$ may come from the same
quadruplet. (The stricter reading ``at most one member per quadruplet''
already fails at $\Q_2$: $101-19=82$, and no single member of $\Q_1$ leaves
a remainder in $\R$.) A member $t$ of $\Q_1$ is written as a royal
expression alone, e.g.\ $11=(2\times3)+5$.
\end{definition}

The following congruences are elementary but shape everything below. All
first members satisfy $p\equiv 11\pmod{30}$, so consecutive first members
differ by a multiple of $30$, and the gaps of the four members of $\Q_k$
satisfy
\[
D\equiv 22,\;24,\;28,\;0 \pmod{30}
\quad\text{for } t=q_k^{(1)},q_k^{(3)},q_k^{(7)},q_k^{(9)} \text{ respectively.}
\]
Every gap is even while every pool member is odd, so a representation with
$r=0$ has an even number of terms. Pool members are $\equiv 11,13,17,19
\pmod{30}$, hence a two-term representation of the gap of a first member
($D\equiv22$) must add two first members ($11+11\equiv22$).

\paragraph{Three ways.}
The utility described in Section~\ref{sec:evidence} computes, for every
member, one representation of each of three kinds:
\begin{itemize}[nosep]
\item[\textbf{A}] \emph{additive}, as in the definition above, chosen so that
  addition is exhausted before any royal multiplication is used (fewer
  multiplications first, then fewer terms, then larger members first);
\item[\textbf{M}] \emph{multiplicative}: any two distinct members of
  $\Pool_k\cup\Q_0$ may be multiplied, and multiplication is preferred over
  addition (the largest unused member is multiplied by the largest member
  that still fits, then the next, plain addition only when nothing can be
  multiplied);
\item[\textbf{E}] \emph{exponential}: in addition, a member may be raised to
  a member power, powers being preferred over products over plain addition.
  Because the order in which powers are tried changes the outcome, two
  variants are kept: E1 takes the largest base first, E2 the largest power
  value first;
\item[\textbf{M1}] \emph{without the base}: the M search applied to the
  member itself instead of to its gap. No base term is required, the pool
  is every member of $\Pool_k\cup\Q_0$ below the target, and the largest
  products carry the number, e.g.
  $854921 = (427249\cdot 2) + (109\cdot 3) + (17\cdot 5) + 11$;
\item[\textbf{E3}] \emph{powers first}: the largest power $b^e$ below the
  member, with $b$ and $e$ distinct members, then again the largest power
  below what is left as long as it covers at least half of it, followed by
  the remainder in the additive way (members below the target added,
  largest first, closed by a flat royal expression), e.g.
  $854921 = 829^2 + 11^5 + 17^3 + 1489 + 107 + 101 + 19$ and
  $19497221 = 11^7 + 2^{13} + 829 + 827 + 199 + 3$. If no choice of
  members closes the remainder, the chain of powers is shortened by one,
  and finally the next largest first power is tried.
\end{itemize}
In A, M, E1 and E2 every member of $\Pool_k\cup\Q_0$ is used at most once
and the base $m_k$ is always a plain term; M1 and E3 also use every member
at most once, the base and exponent of the E3 power included, but have no
base term.

\section{The conjecture}
\label{sec:conj}

We state the observation in three forms. The first is the one the utility
works with; the second and third are the ones we consider interesting,
because they bound the number of terms.

\begin{conjecture}[Ladhe's Quad Conjecture, existence form]
\label{conj:A}
Every member of every prime quadruplet $\Q_k$ with $k\ge2$ has an additive
representation over the chain: $t=m_k+\sum_i s_i+r$ with distinct
$s_i\in\Pool_k$ and $r\in\R\cup\{0\}$.
\end{conjecture}

\begin{conjecture}[bounded pure form]
\label{conj:B}
There is an absolute constant $K$ such that, apart from finitely many gaps,
every gap $D=t-m_k$ ($k\ge2$) is a sum of at most $K$ distinct members of
$\Pool_k$. The data support $K=8$, with exactly ten exceptional gaps:
\[
D\in\{22,\;82,\;84,\;88,\;90,\;172,\;174,\;178,\;180,\;382\},
\]
each of which needs a royal value. No exception occurs beyond $\Q_{435}$
(first member $3{,}512{,}231$) up to $10^9$.
\end{conjecture}

\begin{conjecture}[bounded form with one flat royal value]
\label{conj:C}
Every gap $D=t-m_k$ ($k\ge2$) is a sum of at most five distinct members of
$\Pool_k$ plus at most one element of $\Rflat$; that is, every member has
an additive representation of length at most six whose royal part, if any,
is flat.
\end{conjecture}

\begin{remark}[Why the existence form alone is weak]
Once a few quadruplets exist, the pool is large and the sums of its subsets
cover every even number in range many times over, so Conjecture~\ref{conj:A}
is very likely to hold for reasons that have little to do with
quadruplets; it would also be vacuous if only finitely many quadruplets
existed. What is not automatic is that a \emph{fixed} number of terms
suffices while the gaps grow: the average gap between consecutive
quadruplets near $x$ grows like $(\log x)^4/C_4$, about $35{,}000$ at
$10^9$, and the pool members that fit under a gap are few ($92$ below the
average gap and $336$ below the largest gap, $395{,}520$, in our data).
Conjectures~\ref{conj:B}
and~\ref{conj:C} would be falsified by a single gap that is not a short
sum; the natural falsification test is a very large gap between
consecutive quadruplets.
\end{remark}

\begin{remark}[The exceptional gaps]
All ten exceptions are small and appear early. The four gaps $82,84,88,90$
belong to $\Q_2$ (and reappear whenever two consecutive quadruplets are
$90$ apart, as at $\Q_3$, $\Q_{14}$, $\Q_{437}$, \dots); at that point the
pool is $\Q_1$ alone, whose subset sums cannot reach $82$. The gaps $22$,
$172$--$180$ and $382$ are coverage holes of $\Q_1\cup\Q_2$: for instance
$382=1871-1489$ needs all of $109,107,103,19,17,13,11$ plus $3$. With the
$11$ royal addition values allowed, only $82,84,88,90$ remain exceptional;
with $\Rflat$ allowed there is no exception. Appendix~\ref{app:exc} lists
the ten gaps with their forced representations.
\end{remark}

\section{Coverage: every integer below the next quadruplet}
\label{sec:cover}

The chain also serves as an additive basis for \emph{all} integers, not
only for the quadruplets, once subtraction is allowed inside the royal
part. This section proves that statement from two computed base cases.

\begin{definition}[Signed royal values]
$\R^{\pm}$ is the set of non-negative values of expressions built from
the members $2,3,5,7$ of $\Q_0$, each used at most once, with $+$, $-$ and
$\times$ placed anywhere. It has $83$ elements between $0$ and $210$,
contains every integer from $0$ to $38$, and its first gaps are
$39, 43, 53, 59, 61, 62$. Examples: $0=5-(3+2)$, $1=3-2$, $4=7-3$,
$33=(5\times7)-2$, $55=5\times((2\times7)-3)$.
\end{definition}

\begin{definition}[Derivation]
A \emph{derivation} of an integer $n\ge0$ over $\Q_0,\dots,\Q_k$ is an
identity $n=s_1+\dots+s_j+r$ with distinct $s_i\in\Q_1\cup\dots\cup\Q_k$
and $r\in\R^{\pm}$ ($j=0$ or $r=0$ allowed). Quadruplet members are only
added; subtraction and multiplication occur only among royal members.
Write $C_k$ for the set of integers that have a derivation over
$\Q_0,\dots,\Q_k$ and $T_k$ for the largest $T$ with $[0,T]\subseteq C_k$.
\end{definition}

By direct computation,
\[
T_0=38,\qquad T_1=156,\qquad T_2=576,\qquad T_3=1356,\qquad T_4=4656,
\]
so with the pool $2,3,5,7,11,13,17,19$ every integer from $0$ to $156$ has
a derivation, in particular every integer up to $101$; for instance
$16=11+5$, $53=19+11+(7\times5)-2$ and $55=19+11+(7\times5)$.

\begin{lemma}
\label{lem:extend}
Let $A$ be a set of non-negative integers containing every integer in
$[0,T]$, and let $0<m\le T+1$. Then $A\cup(A+m)$ contains every integer in
$[0,T+m]$.
\end{lemma}
\begin{proof}
An integer $x\in[0,T]$ lies in $A$. An integer $x\in[T+1,T+m]$ satisfies
$0\le x-m\le T$, so $x-m\in A$ and $x\in A+m$.
\end{proof}

\begin{proposition}
\label{prop:cover}
Put $S_3=T_3=1356$ and $S_k=S_{k-1}+4p_k+16$ for $k\ge4$. If
$p_j\le S_{j-1}+1$ for every $j$ with $4\le j\le k+1$, then
$T_k\ge S_k\ge p_{k+1}$: every integer $n$ with $0\le n\le p_{k+1}$ has a
derivation over $\Q_0,\dots,\Q_k$. For $k\le 3$ the same conclusion holds
by the computed values ($T_0\ge p_1=11$, $T_1\ge p_2=101$,
$T_2\ge p_3=191$, $T_3\ge p_4=821$).
\end{proposition}
\begin{proof}
Induction on $k\ge3$ with the claim $T_k\ge S_k$; the case $k=3$ is the
computation. Assume $T_{k-1}\ge S_{k-1}$ and add the members
$p_k,\,p_k+2,\,p_k+6,\,p_k+8$ of $\Q_k$ one at a time, each with
coefficient $0$ or $1$. The first satisfies $p_k\le S_{k-1}+1\le T_{k-1}+1$
by hypothesis, so Lemma~\ref{lem:extend} raises the coverage to
$T_{k-1}+p_k$; each following member is at most $p_k+8$, far below the
coverage reached so far, so the lemma applies three more times and
$T_k\ge T_{k-1}+4p_k+16\ge S_k$. Finally $p_{k+1}\le S_k+1$ and $p_{k+1}$
is odd while $S_k$ is even, hence $p_{k+1}\le S_k$.
\end{proof}

The hypothesis is mild: it fails only if a quadruplet's first member
exceeds $S_{j-1}+1$, which grows like four times the sum of all earlier
first members. In the chain up to $10^9$ it holds at every $j\le 28{,}387$,
and it is tightest at $j=4$ ($p_4=821$ against $S_3+1=1357$). The one-step
condition $p_j\le 5p_{j-1}+16$ is sufficient and also holds throughout
(the largest ratio $p_j/p_{j-1}$ for $j\ge3$ is $4.30$, at $j=4$). So
Proposition~\ref{prop:cover} is unconditional below $10^9$; an
unconditional proof for all $k$ would follow from any Bertrand-type bound
for prime quadruplets, which is open, although far more than what is
needed here.

\begin{remark}
Coverage is a statement about \emph{all} integers with subtraction
allowed and does not imply Conjecture~\ref{conj:A}, which forbids
subtraction, fixes the base $m_k$ and, in its bounded forms, limits the
number of terms. It does show that the differences $D$ of the conjecture,
which are tiny compared with $T_{k-1}$, always have \emph{some} derivation
of this looser kind.
\end{remark}

The utility prints a derivation of any integer with \code{-d}, taking the
largest quad prime below the remainder first, each prime once, and closing
with a flat royal expression (no product containing a sum) whenever the
choice of primes allows it, ranked by fewest multiplications, then fewest
minus signs; a range \code{-d 53 83} prints one line per integer:
\begin{small}\begin{verbatim}
  python3 lc.py -d 16      16 = 13 + 3
  python3 lc.py -d 53      53 = 19 + 17 + 13 + 7 - 3
  python3 lc.py -d 55      55 = 19 + 17 + 13 + 5 + 3 - 2
  python3 lc.py -d 76      76 = 19 + 17 + 13 + 11 + (3*7) - 5
  python3 lc.py -d 99      99 = 19 + 13 + (2*5*7) - 3
  python3 lc.py -d 101     101 = 19 + 17 + 13 + 11 + (5*7) + (2*3)
  python3 lc.py -d 999986171
                           999986171 = 999900529 + 82729 + 2089 + 823 + 3 - 2
\end{verbatim}\end{small}
The greedy choice is not the only derivation; $53=19+11+(7\times5)-2$ and
$55=19+11+(7\times5)$ are equally valid. Up to $20{,}000$ every integer
receives a flat closing.

\section{Worked examples}
\label{sec:examples}

The lines below are copied verbatim from the utility
(\code{python3 lc.py 12 -AME}); \code{A}, \code{M}, \code{M1}, \code{E1},
\code{E2}, \code{E3} label the ways, and a line is omitted when it repeats
an earlier line of the same member (\code{E2} when it coincides with
\code{E1}; in $\Q_1$, which has no base, \code{M1} coincides with
\code{M}). Products are written in parentheses and $b\,\code{\^{}}\,e$
means $b^e$.

\paragraph{$\Q_1=(11,13,17,19)$ from the royal quad.}
The first quadruplet is written with royal expressions only; the utility
also lists alternatives, including ones with subtraction, which the chain
equations of Section~\ref{sec:conj} otherwise never use.
\begin{small}\begin{verbatim}
  11 = (2*3) + 5        also: (2*7) - 3  |  (3*7) - (2*5)  |  7 + 5 + 2 - 3
  13 = (2*5) + 3
  17 = 7 + 5 + 3 + 2
  19 = (2*7) + 5
\end{verbatim}\end{small}

\paragraph{$\Q_2=(101,103,107,109)$: the gap $82$.}
Base $m_2=19$, gap $101-19=82$. The pool is $\Q_1$, whose members other
than the base sum to $41$, so a royal value is unavoidable; $41$ itself is
flat, $(5\times7)+(2\times3)$.
\begin{small}\begin{verbatim}
  A:  101 = 19 + 17 + 13 + 11 + (5*7) + (2*3)
  M:  101 = 19 + (17*3) + (13*2) + 5
  M1: 101 = (19*5) + (3*2)
  E1: 101 = 19 + (7^2) + (11*3)
  E3: 101 = (7^2) + 19 + 17 + 13 + 3
\end{verbatim}\end{small}

\paragraph{$\Q_3=(191,193,197,199)$: the same gap again.}
$191-109=82$, so the representation of the gap is reused with the new
base. The utility stores one equation per distinct gap and reuses it. The
base-free lines M1 and E3 are computed per member and cannot be reused;
here they need only $\Q_1$ and the royal quad.
\begin{small}\begin{verbatim}
  A:  191 = 109 + 17 + 13 + 11 + (5*7) + (2*3)
  M:  191 = 109 + (17*3) + (13*2) + 5
  M1: 191 = (19*7) + (17*3) + 5 + 2
  E1: 191 = 109 + (7^2) + (11*3)
  E3: 191 = (13^2) + (5*3) + 7
\end{verbatim}\end{small}

\paragraph{$\Q_4=(821,823,827,829)$: pure addition.}
Gap $821-199=622=197+193+191+19+17+5$; the three remaining members of
$\Q_3$ are used together with two members of $\Q_1$ and the royal $5$. The
two exponential variants differ here: E1 starts from the largest base
$19^2$, E2 also, but then prefers $3^5=243$ over $5^3=125$.
\begin{small}\begin{verbatim}
  A:  821 = 199 + 197 + 193 + 191 + 19 + 17 + 5
  M:  821 = 199 + (197*3) + (13*2) + 5
  M1: 821 = (199*3) + (103*2) + 13 + 5
  E1: 821 = 199 + (19^2) + (5^3) + 101 + 17 + 11 + 7
  E2: 821 = 199 + (19^2) + (3^5) + 11 + 7
  E3: 821 = (19^2) + (7^3) + 101 + 11 + 5
\end{verbatim}\end{small}

\paragraph{$\Q_6=(1871,\dots)$: a long gap early in the chain.}
Gap $1871-1489=382$ is one of the ten exceptions: no subset of
$\Q_1\cup\Q_2$ sums to $382$, so a royal value is needed. The shortest
representation with a royal \emph{addition} value has seven terms
($199+109+19+17+13+11+14$, with $14=7+5+2$); with one flat royal value
four terms suffice, $382=191+101+17+73$ with $73=(2\times5\times7)+3$.
The utility's addition-first rule prefers the eight-term form below, whose
royal part is the single member $3$. The multiplicative way is one product.
\begin{small}\begin{verbatim}
  A:  1871 = 1489 + 109 + 107 + 103 + 19 + 17 + 13 + 11 + 3
  M:  1871 = 1489 + (191*2)
  M1: 1871 = (829*2) + (19*7) + (17*3) + 13 + 11 + 5
  E1: 1871 = 1489 + (19^2) + (7*3)
  E3: 1871 = (11^3) + 199 + 197 + 109 + 17 + 13 + 5
\end{verbatim}\end{small}

\paragraph{$\Q_{12}=(13001,\dots)$: the typical two-term gap.}
Gap $13001-9439=3562=3461+101$, one member of $\Q_9$ and one of $\Q_2$,
both first members as the congruence $22\equiv 11+11 \pmod{30}$ requires.
About one gap in four has such a two-term representation.
\begin{small}\begin{verbatim}
  A:  13001 = 9439 + 3461 + 101
  M:  13001 = 9439 + (1489*2) + (193*3) + 5
  M1: 13001 = (5659*2) + (199*7) + (19*13) + (5*3) + 17 + 11
  E1: 13001 = 9439 + (19^2) + (13^3) + (109*7) + 107 + 101 + 17 + 11 + 5
  E2: 13001 = 9439 + (13^3) + (19^2) + (109*7) + 107 + 101 + 17 + 11 + 5
  E3: 13001 = (109^2) + 829 + 193 + 19 + 17 + 13 + 11 + (7*5) + 3
\end{verbatim}\end{small}

\paragraph{Two consecutive quadruplets near $3.5\times10^7$.}
$\Q_{2208}$ and $\Q_{2209}$ are $6{,}000$ apart; the gap of the first,
$35545421-35524459=20962$, is $18049+2089+821+3$, and the gap of the second,
$35551421-35545429=5992$, is $5659+199+109+17+5+3$.
\begin{small}\begin{verbatim}
  A:  35545421 = 35524459 + 18049 + 2089 + 821 + 3
  M:  35545421 = 35524459 + (9439*2) + (199*7) + (193*3) + (19*5) + 17
  M1: 35545421 = (17771269*2) + (829*3) + (19*13) + (7*5) + 103 + 11
  E1: 35545421 = 35524459 + (109^2) + (19^3) + (199*7) + (107*5) + 193 + 101
  E3: 35545421 = (5659^2) + (19^5) + (101^3) + 13009 + 1489 + 109 + 103 + 17 + 13

  A:  35551421 = 35545429 + 5659 + 199 + 109 + 17 + 5 + 3
  M:  35551421 = 35545429 + (2089*2) + (199*7) + (103*3) + (19*5) + 17
  M1: 35551421 = (17773369*2) + (1489*3) + (19*11) + 7
  E1: 35551421 = 35545429 + (19^2) + (17^3) + (103*5) + (13*7) + 101 + 11
  E3: 35551421 = (5659^2) + (19^5) + (101^3) + 19423 + 823 + 193 + 191 + 103 + 7
\end{verbatim}\end{small}

The last quadruplet below $10^9$ is $\Q_{28387}=(999986171,\dots)$ with
$999986171 = 999900529 + 82729 + 2089 + 821 + 3$, and the largest gap
between consecutive quadruplets in that range, $395{,}520$ at
$\Q_{19558}$, is still a two-term sum:
$628641701 = 628246189 + 392261 + 3251$.

\section{Computational evidence}
\label{sec:evidence}

\paragraph{The utility.}
\code{lc.py} is a single Python~3 file with no dependencies. It finds
quadruplets with a segmented sieve (only $p\equiv11 \pmod{30}$ is tested),
then searches a representation of every member's gap by iterative
deepening over the number of terms, with the preferences listed in
Section~\ref{sec:defs}, and stores one set of A/M/E1/E2 equations per
distinct gap, plus the base-free M1 and E3 equations of every member, in a
JSON cache that later runs reuse. The additive search is
capped at $16$ terms and $300{,}000$ nodes per call, the M/E searches at
$50{,}000$ nodes, the base-free M1 search, whose pool is the whole chain,
at $1{,}000{,}000$ nodes, and the additive remainder of E3 at $300{,}000$
nodes; a fallback path exists for the case that a search gives up, and
it was never taken in the data reported here (zero fallback flags). Building the chain to $10^9$ took $165$~s on a 2019 Intel laptop
under Python~3.13.5, of which the A/M/E1/E2 equations take $85$~s; the
base-free M1 and E3 equations, computed once per member over the whole
chain, account for the rest.

\paragraph{The checker.}
\code{check.py} re-reads the JSON and, independently of the search code,
verifies every stored equation: the arithmetic, that all members are
distinct, that the base is never reused as a term, and that every term is
smaller than the target. On the $10^9$ dataset it checked $681{,}288$
equations ($6$ per member for all $113{,}548$ members, $\Q_1$ included) and
found $0$ errors.

\paragraph{Minimal lengths.}
\code{stats.py} recomputes, for every distinct gap and independently of the
preferences of the utility, the minimal length $K$ of a representation
under three vocabularies: pool members only (\emph{pure}); pool members
plus at most one value of $\Radd$; pool members plus at most one value of
$\Rflat$. The minimum is taken when the gap is first met, using only the
quadruplets available at that moment. Table~\ref{tab:K} gives the result.

\begin{table}[ht]
\centering
\caption{Minimal length $K$ of a representation of the gap, over the
$13{,}636$ distinct gaps met below $10^9$ (all four members).}
\label{tab:K}
\small
\begin{tabular}{lrrrrrrrrrr}
\toprule
vocabulary & $K{=}1$ & $2$ & $3$ & $4$ & $5$ & $6$ & $7$ & $8$ & none & mean \\
\midrule
pure members            & --   & 2817 & --   & 10375 & --  & 410 & --  & 24 & 10 & 3.65 \\
$+$ one value of $\Radd$  & --   & 2818 & 429  & 9946  & 365 & 53  & 20  & 1  & 4  & 3.59 \\
$+$ one value of $\Rflat$ & 5    & 2927 & 3333 & 7295  & 68  & 8   & --  & -- & 0  & 3.33 \\
\bottomrule
\end{tabular}
\end{table}

The pure column has only even lengths, as it must. The ten gaps without a
pure representation are those of Conjecture~\ref{conj:B}; with $\Radd$
only $82,84,88,90$ remain; with $\Rflat$ nothing remains and six terms
always suffice, which is Conjecture~\ref{conj:C}. The mean minimal length
does not grow along the chain: in eight consecutive blocks of about
$3{,}550$ quadruplets the mean pure $K$ of the first member's gap is
$3.80, 3.75, 3.73, 3.71, 3.71, 3.67, 3.73, 3.71$ and the share of two-term
gaps stays between $22\%$ and $24\%$.

\paragraph{What the utility actually stores.}
Table~\ref{tab:A} summarises the additive representations chosen by the
utility. Because addition is exhausted before any royal product is used,
these are often longer than the minimum: only $120$ of the $113{,}544$
member equations contain a royal product (all with gap $82,84,88$ or
$90$), $52{,}790$ end with an added royal value, and $60{,}634$ are pure
sums of pool members. The longest stored equation has nine terms after
the base. Of the $28{,}386$ first-member gaps, $24{,}977$ repeat a gap seen
earlier in the chain; there are only $13{,}636$ distinct gaps among the
$113{,}544$ member gaps, because consecutive quadruplets are often the
same distance apart.

\begin{table}[ht]
\centering
\caption{Additive representations stored by the utility, $10^9$ dataset.
Length counts the terms after the base, a royal value counting as one.}
\label{tab:A}
\small
\begin{tabular}{lrrrrrrrrr}
\toprule
 & \multicolumn{8}{c}{length} & royal part \\
\cmidrule(lr){2-9}
scope & 2 & 3 & 4 & 5 & 6 & 7 & 8 & 9 & none / $+$ / $\times$ \\
\midrule
first member of each quad & 6493 & -- & 19472 & 30 & 2173 & 119 & 99 & -- & 14776 / 13580 / 30 \\
all four members          & 26875 & 60 & 77222 & 60 & 8323 & 227 & 755 & 22 & 60634 / 52790 / 120 \\
\bottomrule
\end{tabular}
\end{table}

\paragraph{Multiplicative and exponential ways.}
For the $28{,}386$ first-member gaps, the M representations use $90{,}661$
products and $50{,}439$ plain terms, never more than nine terms; $2{,}935$
of the products multiply two royal members. E1 uses $60{,}176$ powers,
$49{,}728$ products and $79{,}246$ plain terms (at most $14$ terms); E2
uses $63{,}213$ powers, $34{,}739$ products and $78{,}861$ plain terms (at
most $12$). E1 and E2 coincide for $9{,}117$ of the $28{,}386$ gaps. These
representations are released for further study; we make no claim about
them beyond their existence within the search caps. The base-free M1 and
powers-first E3 representations were likewise found for every one of the
$113{,}548$ members (only $11$, $17$ and $19$, which are written from the
royal quad alone, admit no power chain and fall back to the power-first
search) and are stored and verified. The E3 chain has one power for
$16{,}926$ members, two for $35{,}775$, three for $46{,}553$ and four for
$14{,}291$; we do not analyse these representations further here.

\section{What the conjecture says about the next quadruplet}
\label{sec:next}

It is tempting to read Conjecture~\ref{conj:B} as a recipe: take the
largest member of the last known quadruplet, add a short sum of earlier
members, and the next quadruplet appears. We want to be precise about what
is and is not true here.

\paragraph{What it gives.}
Every quadruplet after the first carries a \emph{certificate} over its
predecessors: a short equation whose terms are earlier chain members, with
the base term identifying the previous quadruplet. The chain is therefore
self-referential in a strong sense: nothing outside $\Q_0,\dots,\Q_{k-1}$
is needed to write down $\Q_k$ once it is known. The certificate also
turns the search for the next quadruplet into a bounded additive search:
if the conjecture holds with $K$ terms, then the first member of $\Q_k$ is
of the form $m_k+D$ with $D\equiv 22\pmod{30}$ and $D$ a sum of at most
$K$ distinct pool members, and the candidate can be tested by four
primality tests. The utility does exactly this in reverse, producing the
certificate after the sieve has found the quadruplet.

\paragraph{What it does not give.}
The candidates produced by the additive search are far more numerous than
quadruplets: near $10^9$ the number of two-term sums of pool members below
a typical gap is already in the thousands, and almost every admissible $D$
has some short representation. So the conjecture does not point at the
next quadruplet, it only says where it will be found once it exists, and
in practice a sieve remains the faster way to find it. Nor does the
conjecture imply that a next quadruplet exists; infinitude of prime
quadruplets is a separate, open problem.

\paragraph{Why representations are plentiful.}
The Bateman--Horn heuristic predicts roughly $C_4\,y/(\log y)^4$
quadruplets below $y$, and so four times as many pool members. The number
of $j$-element subsets of pool members below $D$ grows like
\[
\frac{1}{j!}\left(\frac{4C_4\,D}{(\log D)^4}\right)^{j},
\]
while their sums range over $O(D)$ even numbers, so for $j\ge2$ the
expected number of $j$-term representations of a given gap grows without
bound as $D$ grows, and grows faster the larger $j$ is. This is the reason the mean minimal length stays flat, and it
suggests the answer to the following question is yes.

\begin{question}
Is every sufficiently large gap a sum of \emph{two} distinct earlier
members? Equivalently, do the lengths $4,6,8$ in Table~\ref{tab:K}
eventually disappear?
\end{question}

\begin{question}
Is a royal value ever needed again beyond the ten exceptional gaps, i.e.\
is Conjecture~\ref{conj:B} true with exactly these exceptions and $K=8$,
or even $K=6$ beyond some point?
\end{question}

\begin{question}
The same construction applies to any prime constellation, for instance
twin primes or sexy prime triplets, with the corresponding admissible
residues. Which constellations admit a bounded chain, and with what $K$?
\end{question}

\begin{question}
The exceptional gaps are exactly the coverage holes of $\Q_1\cup\Q_2$. Is
there a simple description of the set of gaps that are not sums of pool
members at the moment they are first met, for the general constellation?
\end{question}

\section{Availability and reproducibility}
\label{sec:avail}

The utility, the checker and the statistics script are released in the
folder \code{LC} of the repository
\begin{center}\url{https://github.com/SPAlgorithm/LE}\end{center}
The commands below rebuild everything reported here from scratch.
\begin{small}\begin{verbatim}
  python3 lc.py --upto 999999999 --cache paper/qarray_1e9.json   # 85 s
  python3 paper/check.py paper/qarray_1e9.json                   # 0 bad
  python3 paper/stats.py paper/qarray_1e9.json                   # all figures
  python3 lc.py 25 -AME --cache paper/qarray_1e9.json            # the examples
\end{verbatim}\end{small}
The dataset \code{qarray\_1e9.json} ($48{,}068{,}688$ bytes) has SHA-256
\begin{center}\footnotesize\code{72dabebf238cfb792ccb4fc51d126128b497967d9e9b29d68077e7e0cd360477}\end{center}
and \code{quads\_1e9.csv} lists, for each of the $113{,}548$ members
(including $\Q_1$), the base, the gap, the four A/M/E1/E2 equations (the base-free M1/E3
equations are in the JSON only) and the three
minimal lengths. Interactive use: \code{python3 lc.py} asks for a number of
quadruplets; \code{-M}, \code{-E}, \code{-AME} select the ways;
\code{--all} shows all four members; \code{-v} shows which quadruplet each
term comes from; \code{-d N} prints a derivation of any integer $N$ as in
Section~\ref{sec:cover}, extending the chain first if $N$ lies beyond the
cache.

\paragraph{Videos.}
Two narrated explanations accompany this note: \emph{Ladhe's Quad
Conjecture: A New Pattern in Prime Quadruplets} (28~min,
\url{https://youtu.be/_XSOJ0Yj77Q}), which covers Sections
\ref{sec:conj}--\ref{sec:next}, and \emph{Build Any Number From Prime
Quadruplets} (11~min, \url{https://youtu.be/JNWvt3hS540}), which walks
through the coverage proposition of Section~\ref{sec:cover} by hand.

\paragraph{App.}
An interactive companion to this note is being added to the authors' free
iOS app \emph{LE-Games}: \url{https://apps.apple.com/app/id6767880385}.

\paragraph{Conflict of interest.}
The authors declare no competing interests related to this note.

\begin{thebibliography}{9}

\bibitem{HardyLittlewood1923}
G.~H. Hardy and J.~E. Littlewood,
\emph{Some problems of `Partitio numerorum'; III: On the expression of a
number as a sum of primes},
Acta Mathematica 44 (1923), 1--70.

\bibitem{BatemanHorn1962}
P.~T. Bateman and R.~A. Horn,
\emph{A heuristic asymptotic formula concerning the distribution of prime
numbers},
Mathematics of Computation 16 (1962), no.~79, 363--367.

\bibitem{OEIS_A007530}
OEIS Foundation Inc.,
\emph{Sequence A007530: Prime quadruples: numbers $k$ such that $k$, $k+2$,
$k+6$, $k+8$ are all prime},
The On-Line Encyclopedia of Integer Sequences,
\url{https://oeis.org/A007530}.

\bibitem{Zhang2014}
Y.~Zhang,
\emph{Bounded gaps between primes},
Annals of Mathematics 179 (2014), no.~3, 1121--1174.

\bibitem{Maynard2015}
J.~Maynard,
\emph{Small gaps between primes},
Annals of Mathematics 181 (2015), no.~1, 383--413.

\bibitem{Forbes1999}
T.~Forbes,
\emph{Prime clusters and Cunningham chains},
Mathematics of Computation 68 (1999), no.~228, 1739--1747.

\bibitem{LadheConj}
S.~Ladhe and P.~Ladhe,
\emph{Ladhe's Conjecture: additive decompositions of primes},
Zenodo, \url{https://doi.org/10.5281/zenodo.19354450}.

\bibitem{LadheScheme}
S.~Ladhe and P.~Ladhe,
\emph{Ladhe: A compact hash-based one-time signature from additive prime
decompositions},
Zenodo, \url{https://doi.org/10.5281/zenodo.20575084}.

\end{thebibliography}

\appendix

\section{The first 25 quadruplets}
\label{app:table}

For each member: the quad index $k$, the member $t$, the base $m_k$, the gap
$D$, the additive equation stored by the utility (base first), its length
including the base, the minimal pure length $K$ of the gap (``--'' for the
exceptional gaps), and, when the gap was met before, the quad where its
equation was first found. $\Q_1$ is written from royal members only.

\begin{footnotesize}
\begin{longtable}{rrrrlrrr}
\toprule
$k$ & $t$ & $m_k$ & $D$ & additive equation & len & $K$ & from \\
\midrule
\endfirsthead
\toprule
$k$ & $t$ & $m_k$ & $D$ & additive equation & len & $K$ & from \\
\midrule
\endhead
\bottomrule
\endfoot
1 & 11 &  &  & \texttt{(2*3) + 5} & 2 &  &  \\
1 & 13 &  &  & \texttt{(2*5) + 3} & 2 &  &  \\
1 & 17 &  &  & \texttt{7 + 5 + 3 + 2} & 4 &  &  \\
1 & 19 &  &  & \texttt{(2*7) + 5} & 2 &  &  \\
2 & 101 & 19 & 82 & \texttt{19 + 17 + 13 + 11 + (5*7) + (2*3)} & 6 & -- &  \\
2 & 103 & 101 & 2 & \texttt{101 + 2} & 2 &  &  \\
2 & 107 & 101 & 6 & \texttt{101 + (2*3)} & 2 &  &  \\
2 & 109 & 107 & 2 & \texttt{107 + 2} & 2 &  &  \\
3 & 191 & 109 & 82 & \texttt{109 + 19 + 17 + 11 + (5*7)} & 5 & -- & 2 \\
3 & 193 & 191 & 2 & \texttt{191 + 2} & 2 &  &  \\
3 & 197 & 191 & 6 & \texttt{191 + (2*3)} & 2 &  &  \\
3 & 199 & 197 & 2 & \texttt{197 + 2} & 2 &  &  \\
4 & 821 & 199 & 622 & \texttt{199 + 197 + 193 + 191 + 19 + 17 + 5} & 7 & 6 &  \\
4 & 823 & 821 & 2 & \texttt{821 + 2} & 2 &  &  \\
4 & 827 & 821 & 6 & \texttt{821 + (2*3)} & 2 &  &  \\
4 & 829 & 827 & 2 & \texttt{827 + 2} & 2 &  &  \\
5 & 1481 & 829 & 652 & \texttt{829 + 199 + 197 + 193 + 19 + 17 + 13 + 11 + 3} & 9 & 8 &  \\
5 & 1483 & 1481 & 2 & \texttt{1481 + 2} & 2 &  &  \\
5 & 1487 & 1481 & 6 & \texttt{1481 + (2*3)} & 2 &  &  \\
5 & 1489 & 1487 & 2 & \texttt{1487 + 2} & 2 &  &  \\
6 & 1871 & 1489 & 382 & \texttt{1489 + 109 + 107 + 103 + 19 + 17 + 13 + 11 + 3} & 9 & -- &  \\
6 & 1873 & 1871 & 2 & \texttt{1871 + 2} & 2 &  &  \\
6 & 1877 & 1871 & 6 & \texttt{1871 + (2*3)} & 2 &  &  \\
6 & 1879 & 1877 & 2 & \texttt{1877 + 2} & 2 &  &  \\
7 & 2081 & 1879 & 202 & \texttt{1879 + 199 + 3} & 3 & 2 &  \\
7 & 2083 & 2081 & 2 & \texttt{2081 + 2} & 2 &  &  \\
7 & 2087 & 2081 & 6 & \texttt{2081 + (2*3)} & 2 &  &  \\
7 & 2089 & 2087 & 2 & \texttt{2087 + 2} & 2 &  &  \\
8 & 3251 & 2089 & 1162 & \texttt{2089 + 829 + 199 + 109 + 17 + 5 + 3} & 7 & 6 &  \\
8 & 3253 & 3251 & 2 & \texttt{3251 + 2} & 2 &  &  \\
8 & 3257 & 3251 & 6 & \texttt{3251 + (2*3)} & 2 &  &  \\
8 & 3259 & 3257 & 2 & \texttt{3257 + 2} & 2 &  &  \\
9 & 3461 & 3259 & 202 & \texttt{3259 + 199 + 3} & 3 & 2 & 7 \\
9 & 3463 & 3461 & 2 & \texttt{3461 + 2} & 2 &  &  \\
9 & 3467 & 3461 & 6 & \texttt{3461 + (2*3)} & 2 &  &  \\
9 & 3469 & 3467 & 2 & \texttt{3467 + 2} & 2 &  &  \\
10 & 5651 & 3469 & 2182 & \texttt{3469 + 2081 + 101} & 3 & 2 &  \\
10 & 5653 & 5651 & 2 & \texttt{5651 + 2} & 2 &  &  \\
10 & 5657 & 5651 & 6 & \texttt{5651 + (2*3)} & 2 &  &  \\
10 & 5659 & 5657 & 2 & \texttt{5657 + 2} & 2 &  &  \\
11 & 9431 & 5659 & 3772 & \texttt{5659 + 3469 + 199 + 101 + 3} & 5 & 4 &  \\
11 & 9433 & 9431 & 2 & \texttt{9431 + 2} & 2 &  &  \\
11 & 9437 & 9431 & 6 & \texttt{9431 + (2*3)} & 2 &  &  \\
11 & 9439 & 9437 & 2 & \texttt{9437 + 2} & 2 &  &  \\
12 & 13001 & 9439 & 3562 & \texttt{9439 + 3461 + 101} & 3 & 2 &  \\
12 & 13003 & 13001 & 2 & \texttt{13001 + 2} & 2 &  &  \\
12 & 13007 & 13001 & 6 & \texttt{13001 + (2*3)} & 2 &  &  \\
12 & 13009 & 13007 & 2 & \texttt{13007 + 2} & 2 &  &  \\
13 & 15641 & 13009 & 2632 & \texttt{13009 + 1489 + 829 + 199 + 107 + 5 + 3} & 7 & 6 &  \\
13 & 15643 & 15641 & 2 & \texttt{15641 + 2} & 2 &  &  \\
13 & 15647 & 15641 & 6 & \texttt{15641 + (2*3)} & 2 &  &  \\
13 & 15649 & 15647 & 2 & \texttt{15647 + 2} & 2 &  &  \\
14 & 15731 & 15649 & 82 & \texttt{15649 + 19 + 17 + 11 + (5*7)} & 5 & -- & 2 \\
14 & 15733 & 15731 & 2 & \texttt{15731 + 2} & 2 &  &  \\
14 & 15737 & 15731 & 6 & \texttt{15731 + (2*3)} & 2 &  &  \\
14 & 15739 & 15737 & 2 & \texttt{15737 + 2} & 2 &  &  \\
15 & 16061 & 15739 & 322 & \texttt{15739 + 199 + 109 + 11 + 3} & 5 & 4 &  \\
15 & 16063 & 16061 & 2 & \texttt{16061 + 2} & 2 &  &  \\
15 & 16067 & 16061 & 6 & \texttt{16061 + (2*3)} & 2 &  &  \\
15 & 16069 & 16067 & 2 & \texttt{16067 + 2} & 2 &  &  \\
16 & 18041 & 16069 & 1972 & \texttt{16069 + 1871 + 101} & 3 & 2 &  \\
16 & 18043 & 18041 & 2 & \texttt{18041 + 2} & 2 &  &  \\
16 & 18047 & 18041 & 6 & \texttt{18041 + (2*3)} & 2 &  &  \\
16 & 18049 & 18047 & 2 & \texttt{18047 + 2} & 2 &  &  \\
17 & 18911 & 18049 & 862 & \texttt{18049 + 829 + 19 + 11 + 3} & 5 & 4 &  \\
17 & 18913 & 18911 & 2 & \texttt{18911 + 2} & 2 &  &  \\
17 & 18917 & 18911 & 6 & \texttt{18911 + (2*3)} & 2 &  &  \\
17 & 18919 & 18917 & 2 & \texttt{18917 + 2} & 2 &  &  \\
18 & 19421 & 18919 & 502 & \texttt{18919 + 199 + 197 + 103 + 3} & 5 & 4 &  \\
18 & 19423 & 19421 & 2 & \texttt{19421 + 2} & 2 &  &  \\
18 & 19427 & 19421 & 6 & \texttt{19421 + (2*3)} & 2 &  &  \\
18 & 19429 & 19427 & 2 & \texttt{19427 + 2} & 2 &  &  \\
19 & 21011 & 19429 & 1582 & \texttt{19429 + 1481 + 101} & 3 & 2 &  \\
19 & 21013 & 21011 & 2 & \texttt{21011 + 2} & 2 &  &  \\
19 & 21017 & 21011 & 6 & \texttt{21011 + (2*3)} & 2 &  &  \\
19 & 21019 & 21017 & 2 & \texttt{21017 + 2} & 2 &  &  \\
20 & 22271 & 21019 & 1252 & \texttt{21019 + 829 + 199 + 197 + 19 + 5 + 3} & 7 & 6 &  \\
20 & 22273 & 22271 & 2 & \texttt{22271 + 2} & 2 &  &  \\
20 & 22277 & 22271 & 6 & \texttt{22271 + (2*3)} & 2 &  &  \\
20 & 22279 & 22277 & 2 & \texttt{22277 + 2} & 2 &  &  \\
21 & 25301 & 22279 & 3022 & \texttt{22279 + 2089 + 829 + 101 + 3} & 5 & 4 &  \\
21 & 25303 & 25301 & 2 & \texttt{25301 + 2} & 2 &  &  \\
21 & 25307 & 25301 & 6 & \texttt{25301 + (2*3)} & 2 &  &  \\
21 & 25309 & 25307 & 2 & \texttt{25307 + 2} & 2 &  &  \\
22 & 31721 & 25309 & 6412 & \texttt{25309 + 3257 + 1483 + 1481 + 191} & 5 & 4 &  \\
22 & 31723 & 31721 & 2 & \texttt{31721 + 2} & 2 &  &  \\
22 & 31727 & 31721 & 6 & \texttt{31721 + (2*3)} & 2 &  &  \\
22 & 31729 & 31727 & 2 & \texttt{31727 + 2} & 2 &  &  \\
23 & 34841 & 31729 & 3112 & \texttt{31729 + 2089 + 829 + 191 + 3} & 5 & 4 &  \\
23 & 34843 & 34841 & 2 & \texttt{34841 + 2} & 2 &  &  \\
23 & 34847 & 34841 & 6 & \texttt{34841 + (2*3)} & 2 &  &  \\
23 & 34849 & 34847 & 2 & \texttt{34847 + 2} & 2 &  &  \\
24 & 43781 & 34849 & 8932 & \texttt{34849 + 5659 + 3259 + 11 + 3} & 5 & 4 &  \\
24 & 43783 & 43781 & 2 & \texttt{43781 + 2} & 2 &  &  \\
24 & 43787 & 43781 & 6 & \texttt{43781 + (2*3)} & 2 &  &  \\
24 & 43789 & 43787 & 2 & \texttt{43787 + 2} & 2 &  &  \\
25 & 51341 & 43789 & 7552 & \texttt{43789 + 5659 + 1879 + 11 + 3} & 5 & 4 &  \\
25 & 51343 & 51341 & 2 & \texttt{51341 + 2} & 2 &  &  \\
25 & 51347 & 51341 & 6 & \texttt{51341 + (2*3)} & 2 &  &  \\
25 & 51349 & 51347 & 2 & \texttt{51347 + 2} & 2 &  &  \\

\end{longtable}
\end{footnotesize}

\section{The ten exceptional gaps}
\label{app:exc}

Gaps that are not a sum of distinct pool members at the moment they are
first met, with the additive (A), multiplicative (M) and exponential (E1)
representations of the gap stored by the utility.

\begin{scriptsize}
\setlength{\tabcolsep}{3pt}
\begin{tabular}{rrrlll}
\toprule
$D$ & $k$ & $t$ & A & M & E1 \\
\midrule
22  & 169 & 1006331 & \code{19 + 3} & \code{(11*2)} & \code{(3\^{}2) + 13} \\
82  & 2   & 101     & \code{17 + 13 + 11 + (5*7) + (2*3)} & \code{(17*3) + (13*2) + 5} & \code{(7\^{}2) + (11*3)} \\
84  & 2   & 103     & \code{11 + (2*5*7) + 3} & \code{(17*3) + (13*2) + 7} & \code{(7\^{}2) + 17 + 13 + 5} \\
88  & 2   & 107     & \code{17 + 13 + 11 + (2*3*7) + 5} & \code{(17*5) + 3} & \code{(7\^{}2) + (13*3)} \\
90  & 2   & 109     & \code{17 + (2*5*7) + 3} & \code{(17*5) + 3 + 2} & \code{(7\^{}2) + 17 + 13 + 11} \\
172 & 435 & 3512231 & \code{109 + 19 + 17 + 13 + 11 + 3} & \code{(19*7) + (17*2) + 5} & \code{(13\^{}2) + 3} \\
174 & 435 & 3512233 & \code{109 + 19 + 17 + 13 + 11 + 5} & \code{(19*7) + (13*3) + 2} & \code{(13\^{}2) + 5} \\
178 & 435 & 3512237 & \code{109 + 19 + 17 + 13 + 11 + 7 + 2} & \code{(19*7) + (17*2) + 11} & \code{(11\^{}2) + (19*3)} \\
180 & 435 & 3512239 & \code{103 + 19 + 17 + 13 + 11 + 7 + 5 + 3 + 2} & \code{(19*7) + (17*2) + 13} & \code{(13\^{}2) + 11} \\
382 & 6   & 1871    & \code{109 + 107 + 103 + 19 + 17 + 13 + 11 + 3} & \code{(191*2)} & \code{(19\^{}2) + (7*3)} \\
\bottomrule
\end{tabular}
\end{scriptsize}

\section{Notation}

\begin{small}
\begin{tabular}{ll}
\toprule
$\Q_0$ & the royal quad $(2,3,5,7)$ \\
$\Q_k$, $k\ge1$ & the $k$-th prime quadruplet $(p,p+2,p+6,p+8)$, $p\ge11$ \\
$q_k^{(d)}$ & the member of $\Q_k$ ending in digit $d\in\{1,3,7,9\}$ \\
$m_k=q_{k-1}^{(9)}$ & the base of $\Q_k$ \\
$D=t-m_k$ & the gap of a member $t$ of $\Q_k$ \\
$\Pool_k=\Q_1\cup\dots\cup\Q_{k-1}$ & the pool available to $\Q_k$ \\
$\R$, $\Rflat$, $\Radd$ & royal values: all ($68$), flat ($38$), additive ($11$) \\
$\R^{\pm}$ & non-negative royal values with $+,-,\times$ ($83$) \\
$K$ & minimal length of a representation of a gap \\
$C_k$, $T_k$ & integers with a derivation over $\Q_0..\Q_k$; largest $T$ with $[0,T]\subseteq C_k$ \\
$S_k$ & $S_3=T_3$, $S_k=S_{k-1}+4p_k+16$: the proven lower bound for $T_k$ \\
\bottomrule
\end{tabular}
\end{small}

\end{document}
